\subsection{Recurrent Importance Sampling} \label{sup:recpolgrad}

In off-policy reinforcement learning (RL), the replay buffer is populated with trajectories generated by past policies, which may differ from the current policy. A central question is whether such transitions can still be used to update the current policy. This is addressed via \emph{importance sampling}, which adjusts for the distribution mismatch between the behavior and target policies.

The policy gradient is defined as:
\begin{equation}
    \nabla_{\theta} J(\pi_{\theta}) = \nabla_{\theta} \mathbb{E}_{\tau \sim \pi_{\theta}} \left[ R(\tau) \right]
    = \nabla_{\theta} \int_{\tau} P(\tau \mid \theta) R(\tau)
\end{equation}
where \( \tau = (s_0, a_0, s_1, a_1, \ldots, s_T, a_T) \) is a trajectory, and \( P(\tau \mid \theta) \) denotes the probability of observing \( \tau \) under the current policy.

Expanding this probability for a Markovian environment:
\begin{equation} \label{eq:transition}
    P(\tau \mid \theta) = \rho_0(s_0) \prod_{t=0}^{T} P(s_{t+1} \mid s_t, a_t)\, \pi_{\theta}(a_t \mid s_t)
\end{equation}
where \( \rho_0 \) is the initial state distribution, \( P \) is the environment dynamics, and \( \pi_{\theta} \) is the policy, parametrized by parameters \( \theta \).

Importance sampling allows us to rewrite expectations under \( \pi_{\theta} \) using samples from a different policy \( \pi_{\theta'} \):
\begin{equation} \label{eq:importance}
    \mathbb{E}_{\tau \sim \pi_{\theta}} \left[ R(\tau) \right] =
    \mathbb{E}_{\tau \sim \pi_{\theta'}} \left[
        \frac{P(\tau \mid \theta)}{P(\tau \mid \theta')} R(\tau)
    \right]
\end{equation}
This forms the basis for off-policy policy gradient methods.

\paragraph{Recurrent Policies.}
In the recurrent (or spiking) setting, policies depend on a history of observations or states. We can denote a recurrent policy as \( \pi_{\theta}(a_t \mid h_t) \), where \( h_t = f(s_0, \ldots, s_t) \) summarizes past information, such as via hidden states in RNNs or membrane potential in SNNs.

Replacing \( \pi_{\theta}(a_t \mid s_t) \) with \( \pi_{\theta}(a_t \mid h_t) \) in \eqref{eq:transition}, we obtain:
\begin{equation}
    P(\tau \mid \theta) = \rho_0(s_0) \prod_{t=0}^{T} P(s_{t+1} \mid h_t, a_t)\, \pi_{\theta}(a_t \mid h_t)
\end{equation}
which we can substitute into \eqref{eq:importance}:
\begin{equation} \label{eq:recurrent-importance}
    \mathbb{E}_{\tau \sim \pi_{\theta}} \left[ R(\tau) \right] =
    \mathbb{E}_{\tau \sim \pi_{\theta'}} \left[
        \prod_{t=0}^T \frac{\pi_{\theta}(a_t \mid h_t)}{\pi_{\theta'}(a_t \mid h_t)} R(\tau)
    \right]
\end{equation}

\paragraph{Conclusion.}
Equation \eqref{eq:recurrent-importance} shows that the standard policy gradient theorem remains valid for recurrent or spiking policies. The only modification is that the policy’s input includes temporal context. As long as this is accounted for in the importance weights, off-policy learning remains theoretically sound. This justifies the use of off-policy experience replay in recurrent or spiking RL agents. 